
\section{Conclusion and Future Work}
\label{sec:conclusion}

This paper presented a reproducible research framework
for evaluating the impact of Generative AI on professional
writing productivity and work quality. We proposed a
randomized, counterbalanced crossover design, a fair
no-AI baseline, and predefined metrics for completion
time, output quality, and factual correctness.

Regarding the primary research question, the synthetic
demonstration produced a mean paired time saving of
3.875 minutes, with a 95\% bootstrap confidence interval
of [0.00, 6.25] minutes. However, these values were
generated from four fictional participant pairs and
cannot establish a real-world productivity improvement.
The secondary question concerning quality preservation
also remains unresolved without real participant data
and formal non-inferiority testing.

Future work will focus on three directions. First,
we will recruit consenting participants, pilot-test
the equivalence of writing tasks, and determine an
appropriate sample size. Second, we will conduct
blinded quality assessments and evaluate factual
Precision, Recall, and F1-score using real responses.
Third, we will investigate task-order effects and
extend the protocol to additional professional
writing scenarios.

Overall, the main contribution of this work is a
transparent experimental protocol and a tested
Python analysis pipeline. These components provide
a foundation for a future empirical investigation
rather than evidence that Generative AI already
improves employee productivity.
